\section{Introduction}
\label{sec:introduction}

Finite-window observations provide a common interface for symbolic dynamics, automata, coding theory, and finite-resolution data acquisition.  Given a phase space \(X\), a singly generated action
\[
T^t:X\to X,
\]
and a measurable binary window \(W\subseteq X\), the readout
\[
s_t(x)=\one_W(T^t x)
\]
produces, for each resolution \(m\), a finite word
\[
\omega_m(x;t)=s_t(x)s_{t+1}(x)\cdots s_{t+m-1}(x)\in\Om_m,
\qquad
\Om_m=\{0,1\}^m.
\]
The difficulty is immediate:
\[
|\Om_m|=2^m.
\]
Thus the observable microstate space grows exponentially with resolution.  A direct identification of microstates with stable types is unstable: most microstates become sparse, cross-scale comparisons are not canonical, and finite-window evidence is difficult to audit across resolutions.

A natural response is to replace raw microstates by canonical stable representatives.  This requires a map
\[
\Fold_m:\Om_m\to X_m
\]
from microstates to stable type spaces, together with a normal-form theorem and a cross-resolution compatibility rule.  Such a folding map should not be an arbitrary clustering procedure.  It should be computable, idempotent, canonical, and compatible with the finite observable structure generated by the scan.

This paper uses the golden-mean/Fibonacci syntax as a canonical normalization grammar.  Equivalently, the stable representatives are organized by Zeckendorf legality: each residue class modulo the finite Fibonacci modulus has a unique legal representative.  The golden mean is not used as a physical postulate.  It is used as a canonical anti-resonant baseline and as a finite normalization grammar whose local rewriting implementation can be audited.

\paragraph{Main contributions.}
The paper makes five contributions.
\begin{enumerate}[label=(\roman*)]
\item We construct golden-mean folding maps \(\Fold_m:\Om_m\to X_m\) as finite Zeckendorf residue sections, and prove well-definedness, computability, idempotence, surjectivity, and uniqueness of stable representatives.
\item We construct fold-aware restriction maps by normalizing first and restricting second, and prove that the folded type spaces are compatible across adjacent resolutions.
\item We formulate recursive addressing as an internal re-selection and organization of previous readout layers, and prove that derived addresses do not enlarge the cumulative visible sigma-algebra.
\item We define the scan error profile \(\varepsilon_m(P;\mu)\), prove its cylinder decomposition, identify the Bayes-optimal prefix classifier, and give boundary-dimension clarity bounds.
\item We provide finite graph, entropy, spectral-fingerprint, and reproducible-experiment interfaces so that the stable-type generator can be independently audited.
\end{enumerate}

\paragraph{What is not claimed.}
The claims of this paper are deliberately restricted.  We do not claim that the golden mean by itself yields quantum mechanics, spacetime, or Einstein equations.  We also do not use Hilbert-space, quantum, or spacetime assumptions in the proofs of the main theorems.  The paper establishes a mathematical and operational front end: binary readouts can be folded into stable types with normal forms, cross-scale compatibility, recursive addressing, and statistical clarity certificates.  Downstream arithmetic, spectral, \v{C}ech-residual, Hilbert, or spacetime bridges require additional admissibility conditions and are left outside the main theorem chain.

\paragraph{Organization.}
\Cref{sec:scope} fixes the scope and claim hierarchy.  \Cref{sec:scan-projection} defines scan--projection systems, window microstates, cylinder events, and visible sigma-algebras.  \Cref{sec:golden-syntax} recalls the golden-mean and Zeckendorf syntax.  \Cref{sec:folding} defines the folding map and proves the normal-form theorem.  \Cref{sec:cross-resolution} establishes cross-resolution compatibility through fold-aware restriction.  \Cref{sec:recursive-addressing} develops recursive addressing and proves visible sigma-algebra non-expansion.  \Cref{sec:clarity} defines scan error profiles and proves clarity bounds.  \Cref{sec:certificates} gives finite graph and certificate interfaces.  \Cref{sec:experiments} states reproducible experiment protocols.  \Cref{sec:discussion} discusses limitations and downstream bridges, and \Cref{sec:conclusion} summarizes the main claim.
